% Einheit 3 von 3 – Quadratwurzeln (bank/potenzen-wurzeln/e3.jsonl, zusammenbau v0.1)
%% TODO Zweigzeile Teil 1: was hier gelernt wird (Satz in Schülersprache aus dem Einheitstitel) – Einheit 3
\einheitenkopf[e3]{Einheit 3 von 3 · Quadratwurzeln}
\zweigzeile{ab Kl. 7, je nach Buch bis Kl. 9 (am Gymnasium ab Kl. 8) · P10}
\setcounter{aufgabe}{20}

%% TODO Titel: Ich-kann-Satz fehlt in der Bank (Platzhalter: Kettenname) – Nr. 21
%% TODO Anweisung: ein Satz über der ersten Teilaufgabe fehlt in der Bank – Nr. 21
\begin{aufgabe}{Was zuerst?}
\begin{teile}
\teil $\sqrt{(-7)^2}$ – was zuerst?\\ \kreuz{erst die Wurzel} \\ \kreuz{erst das Quadrat}
\end{teile}
\end{aufgabe}

%% TODO Titel: Ich-kann-Satz fehlt in der Bank (Platzhalter: Kettenname) – Nr. 22
%% TODO Anweisung: ein Satz über der ersten Teilaufgabe fehlt in der Bank – Nr. 22
\begin{aufgabe}{Quadratwurzeln}
\begin{teile}
\teil $\sqrt{121}$ – geht die Wurzel im Kopf auf?\\ \kreuz{ja, Quadratzahl} \\ \kreuz{nein, Näherungswert}
\end{teile}
%% TODO Darstellung für schwach fehlt in der Bank (potenzen-wurzeln-e3-k2-s1-v1; Abschnitt „Für schwache Schüler“)
\swz{$\sqrt{49}$ – wie viel? Mache die Probe.}{}
%% TODO Darstellung für schwach fehlt in der Bank (potenzen-wurzeln-e3-k2-s1-v2; Abschnitt „Für schwache Schüler“)
\swz{$\sqrt{64}$ – wie viel? Mache die Probe.}{}
%% TODO Darstellung für schwach fehlt in der Bank (potenzen-wurzeln-e3-k2-s1-v3; Abschnitt „Für schwache Schüler“)
\swz{$\sqrt{100}$ – wie viel? Mache die Probe.}{}
%% TODO Darstellung für schwach fehlt in der Bank (potenzen-wurzeln-e3-k2-s1-v4; Abschnitt „Für schwache Schüler“)
\swz{$\sqrt{225}$ – wie viel? Mache die Probe.}{}
%% TODO Darstellung für schwach fehlt in der Bank (potenzen-wurzeln-e3-k2-s1-v5; Abschnitt „Für schwache Schüler“)
\swz{$\sqrt{400}$ – wie viel? Mache die Probe.}{}
\swfrage{Was bleibt gleich, was ändert sich?}
%% TODO Darstellung für schwach fehlt in der Bank (potenzen-wurzeln-e3-k2-s2-v1; Abschnitt „Für schwache Schüler“)
\swz{$\sqrt{11}$ – mit dem Taschenrechner, auf zwei Stellen nach dem Komma gerundet?}{}
%% TODO Darstellung für schwach fehlt in der Bank (potenzen-wurzeln-e3-k2-s3-v1; Abschnitt „Für schwache Schüler“)
\swz{$\sqrt{60}$ – zwischen welchen ganzen Zahlen? Prüfe dann mit dem Taschenrechner.}{}
%% TODO Darstellung für schwach fehlt in der Bank (potenzen-wurzeln-e3-k2-s4-v1; Abschnitt „Für schwache Schüler“)
\swz{$\sqrt{3} \;\square\; 1{,}7$ – $<$, $=$ oder $>$?}{}
%% TODO Darstellung für schwach fehlt in der Bank (potenzen-wurzeln-e3-k2-s5-v1; Abschnitt „Für schwache Schüler“)
\swz{$\sqrt{7}$; $2{,}6$; $\frac{5}{2}$; $-2{,}7$ – aufsteigend geordnet?}{}
%% TODO Darstellung für schwach fehlt in der Bank (potenzen-wurzeln-e3-k2-s6-v1; Abschnitt „Für schwache Schüler“)
\swz{$\left(\sqrt{13}\right)^2$ – wie viel?}{}
%% TODO Darstellung für schwach fehlt in der Bank (potenzen-wurzeln-e3-k2-s7-v1; Abschnitt „Für schwache Schüler“)
\swz{$\sqrt{(-8)^2}$ – wie viel?}{}
%% TODO Darstellung für schwach fehlt in der Bank (potenzen-wurzeln-e3-k2-s8-v1; Abschnitt „Für schwache Schüler“)
\swz{$\sqrt{-49}$ und $-\sqrt{49}$ – welcher Term hat einen Wert? Gib ihn an.}{}
%% TODO Darstellung für schwach fehlt in der Bank (potenzen-wurzeln-e3-k2-s9-v1; Abschnitt „Für schwache Schüler“)
\swz{$\sqrt{0{,}49}$ – wie viel?}{}
%% TODO Darstellung für schwach fehlt in der Bank (potenzen-wurzeln-e3-k2-s10-v1; Abschnitt „Für schwache Schüler“)
\swz{Ein quadratischer Teppich hat einen Flächeninhalt von $6{,}25\,\text{m}^2$. Wie lang ist eine Seite? \hfill (P10 2020 OS)}{}
%% TODO Darstellung für schwach fehlt in der Bank (potenzen-wurzeln-e3-k2-s11-v1; Abschnitt „Für schwache Schüler“)
\swz{$\sqrt{7^2 + 24^2}$ – wie viel?}{}
%% TODO Darstellung für schwach fehlt in der Bank (potenzen-wurzeln-e3-k2-s12-v1; Abschnitt „Für schwache Schüler“)
\swz{$\sqrt[3]{64}$ – wie viel?}{}
\begin{teile}
\teil Welche Aussage ist wahr? Kreuze an. \hfill (P10 2015 OS)\\ \kreuz{$\sqrt{(-7)^2} = -7$} \\ \kreuz{$\sqrt{(-7)^2} = 7$} \\ \kreuz{$\sqrt{(-7)^2}$ ist nicht definiert}
\end{teile}
\end{aufgabe}

%% TODO Titel: Ich-kann-Satz fehlt in der Bank (Platzhalter: Kettenname) – Nr. 23
%% TODO Anweisung: ein Satz über der ersten Teilaufgabe fehlt in der Bank – Nr. 23
\begin{aufgabe}{Quadratwurzeln – Fehler finden, Begründen}
\swz[3]{Emil rechnet $\sqrt{64} = 32$. Wo ist der Fehler?}{}
\swfrage{Begründe, warum $\sqrt{49}$ nicht $-7$ ist, obwohl $(-7)^2 = 49$ gilt.}
\end{aufgabe}

\uebersichtskasten{Wurzel: Die Quadratwurzel ist die Umkehrung des Quadrierens – √a ist die Zahl, die mit sich selbst malgenommen a ergibt, und sie ist nie negativ: √144 = 12, weil 12 · 12 = 144. \\ Quadrat und Wurzel heben sich auf – bei negativer Basis erst das Quadrat rechnen, die Wurzel bleibt positiv: √(9²) = 9 und √((−9)²) = √81 = 9. \\ Keine Wurzel aus einer negativen Zahl: √(−81) hat keinen Wert, weil kein Quadrat negativ ist. \\ Näherungswert: Geht die Wurzel nicht auf, liegt sie zwischen zwei Nachbar-Quadratzahlen; der Taschenrechner gibt den Wert, gerundet wird erst am Ende: √30 ≈ 5,48 (zwischen √25 = 5 und √36 = 6).}
